% New material for proof.tex.
% The bibliography entry Straub2014, printed at the end of this file as a
% comment, should be inserted into the main thebibliography environment.

\section{A diagonal Cartier representation}
\label{sec:cartier-diagonal}

Put
\[
 \mathscr F(x,y,z,w)
  =(1-x-y)(1-z-w)-xyzw.
\]
The following observation places every truncated Ap\'ery zero set inside
one fixed algebraic family.

\begin{theorem}[Diagonal Cartier representation]
\label{thm:cartier-diagonal}
Let \(p\ge5\) be prime.  For every \(0\le r\le p-1\),
\[
 b_r\equiv
 [x^ry^rz^rw^r]\,\mathscr F(x,y,z,w)^{p-1}
 \pmod p.
\]
\end{theorem}

\begin{proof}
Straub's diagonal formula~\cite[Theorem~1.1]{Straub2014} gives
\[
 b_r=[x^ry^rz^rw^r]\,\frac1{\mathscr F(x,y,z,w)}.
\]
In \(\mathbb F_p[[x,y,z,w]]\), the freshman's dream gives
\[
 \mathscr F^{p-1}
 =\frac{\mathscr F^p}{\mathscr F}
 =\frac{\mathscr F(x^p,y^p,z^p,w^p)}{\mathscr F(x,y,z,w)}.
\]
Every nonconstant monomial in the numerator on the right has at least
one exponent at least \(p\).  Since \(0\le r<p\), only its constant
term contributes to \([x^ry^rz^rw^r]\).  The asserted congruence
follows.
\end{proof}

\begin{remark}
The first displayed identity can also be checked directly by expanding
\[
 \frac1{\mathscr F}
 =\sum_{j\ge0}
 \frac{(xyzw)^j}
 {(1-x-y)^{j+1}(1-z-w)^{j+1}}.
\]
Its diagonal coefficient is
\[
 \sum_{j=0}^r
 \left(\frac{(2r-j)!}{j!(r-j)!^2}\right)^2
 =\sum_{k=0}^r
 \binom rk^2\binom{r+k}{k}^2=b_r.
\]
The proof of Theorem~\ref{thm:cartier-diagonal} is the corresponding
one-digit Cartier extraction.
\end{remark}


\section{Fourier non-concentration of the Ap\'ery zero sets}
\label{sec:fourier-nonconcentration}

For a prime \(p\ge7\), write
\[
 \mathcal Z_p=\{0\le r<p:b_r\equiv0\pmod p\},
 \qquad Z(p)=|\mathcal Z_p|,
\]
and define
\[
 F_p(k)=\sum_{r\in\mathcal Z_p}e_p(kr),
 \qquad e_p(t)=e^{2\pi i t/p}.
\]

\begin{lemma}[Fourier non-concentration]
\label{lem:fourier-nonconcentration}
For every prime \(p\ge7\) and every integer \(1\le K\le p\),
\[
 \sum_{1\le |k|\le K}|F_p(k)|^2
 \le
 4\pi^2\left(KZ(p)+\frac{p^2}{K}\right).
\]
Here \(k\) ranges over ordinary integers and \(F_p\) is extended
periodically modulo \(p\).
\end{lemma}

\begin{proof}
For \(1\le h<p\), let
\[
 A_p(h)
 =\#\{r:0\le r<r+h<p,\ r,r+h\in\mathcal Z_p\}.
\]
Lemma~\ref{lem:no-consec} gives \(A_p(1)=0\).  For \(h\ge2\),
every pair counted by \(A_p(h)\) makes \(r\) a root modulo \(p\) of
the nonzero gap polynomial \(N_h\).  Lemmas~\ref{lem:gap-poly} and
\ref{lem:nonvanish} therefore give
\begin{equation}
 A_p(h)\le3(h-1).
 \label{eq:new-fourier-gap}
\end{equation}

Suppose first that \(K\le p/8\), and put
\[
 H=\left\lfloor\frac{p}{4K}\right\rfloor.
\]
Since \(p/(4K)\ge2\),
\begin{equation}
 \frac{p}{8K}\le H\le\frac{p}{4K},
 \qquad \frac{HK}{p}\le\frac14.
 \label{eq:new-fourier-H}
\end{equation}
Split \(\mathcal Z_p=Z_0\sqcup Z_1\), where
\[
 Z_0\subseteq[0,(p-1)/2],
 \qquad
 Z_1\subseteq[(p+1)/2,p-1],
\]
and set \(F_i(k)=\sum_{r\in Z_i}e_p(kr)\).  Then
\begin{equation}
 |F_p(k)|^2\le2|F_0(k)|^2+2|F_1(k)|^2.
 \label{eq:new-fourier-halves}
\end{equation}
This division into half intervals removes cyclic wrap-around: within
either half, a congruence \(s-r\equiv h\pmod p\) with
\(|h|<H\le p/4\) is the ordinary equality \(s-r=h\).

Consider the finite Fej\'er kernel
\[
 \mathcal F_H(k)
 =\frac1H\left|\sum_{u=0}^{H-1}e_p(ku)\right|^2
 =\sum_{|h|<H}
 \left(1-\frac{|h|}{H}\right)e_p(kh).
\]
For \(1\le|k|\le K\), the sine formula and
\eqref{eq:new-fourier-H} imply
\[
 \left|\sum_{u=0}^{H-1}e_p(ku)\right|
 =\frac{|\sin(\pi Hk/p)|}{|\sin(\pi k/p)|}
 \ge\frac{2H}{\pi}.
\]
Indeed, \(\sin(\pi x)\ge2x\) for \(0\le x\le1/4\), while
\(\sin(\pi x)\le\pi x\) for \(x\ge0\).  Consequently
\begin{equation}
 \mathbf1_{\{1\le|k|\le K\}}
 \le\frac{\pi^2}{4H}\mathcal F_H(k).
 \label{eq:new-fourier-majorant}
\end{equation}

For \(i=0,1\), let
\[
 A_i(h)=\#\{r:r,r+h\in Z_i\},
\]
where ordinary integer representatives are used.  Finite Fourier
orthogonality and \eqref{eq:new-fourier-majorant} give
\begin{align*}
 \sum_{1\le|k|\le K}|F_i(k)|^2
 &\le\frac{\pi^2}{4H}
   \sum_{k\bmod p}\mathcal F_H(k)|F_i(k)|^2\\
 &=\frac{\pi^2p}{4H}
   \left(
     |Z_i|+
     2\sum_{h=1}^{H-1}
       \left(1-\frac hH\right)A_i(h)
   \right).
\end{align*}
Summing this inequality over \(i\), using
\eqref{eq:new-fourier-halves}, and observing that
\(A_0(h)+A_1(h)\le A_p(h)\), we obtain
\begin{align*}
 \sum_{1\le|k|\le K}|F_p(k)|^2
 &\le\frac{\pi^2p}{2H}
   \left(Z(p)+2\sum_{h=1}^{H-1}A_p(h)\right)\\
 &\le\frac{\pi^2p}{2H}\left(Z(p)+3H^2\right).
\end{align*}
The last step follows from \eqref{eq:new-fourier-gap}, since
\[
 \sum_{h=2}^{H-1}3(h-1)
 =\frac32(H-1)(H-2)\le\frac32H^2.
\]
Using \eqref{eq:new-fourier-H} now yields
\[
 \sum_{1\le|k|\le K}|F_p(k)|^2
 \le
 4\pi^2KZ(p)+\frac{3\pi^2}{8}\frac{p^2}{K},
\]
which is stronger than the claimed estimate.

It remains to treat \(p/8<K\le p\).  Each residue modulo \(p\)
occurs at most twice among the integers \(k\) with
\(1\le|k|\le K\).  Parseval therefore gives
\[
 \sum_{1\le|k|\le K}|F_p(k)|^2
 \le2\sum_{a\bmod p}|F_p(a)|^2
 =2pZ(p)
 <16KZ(p).
\]
Since \(16<4\pi^2\), this is again bounded by the asserted
right-hand side.
\end{proof}

\begin{remark}
Applying the Fej\'er expansion directly to \(\mathcal Z_p\) would
also count ordinary gaps \(p-h\) at the coefficient indexed by a
small cyclic gap \(h\).  The degree bound
\(A_p(h)\le3(h-1)\) does not control those wrap-around pairs.
The two-half partition in the proof is therefore essential.
\end{remark}


\section{Moment collapse and a strong anti-tail hypothesis}
\label{sec:moment-collapse}

For a dyadic block
\[
 \mathcal P_X=\{p\text{ prime}:X<p\le2X\},
\]
recall
\[
 K_X(m)
 =\#\{p\in\mathcal P_X:m\bmod p\in\mathcal Z_p\},
 \qquad
 \lambda_X=\sum_{p\in\mathcal P_X}\frac{Z(p)}p,
\]
and put
\[
 M_X=\max_{0\le m<X^2}K_X(m).
\]
For an integer \(u\ge0\), write
\((u)_j=u(u-1)\cdots(u-j+1)\).

\begin{proposition}[Moment collapse]
\label{prop:moment-collapse}
For every fixed integer \(j\ge2\),
\[
 \sum_{0\le m<X^2}(K_X(m))_j
 \le
 M_X^{j-2}
 \sum_{0\le m<X^2}(K_X(m))_2
 \le
 5X^2\lambda_X^2M_X^{j-2}.
\]
\end{proposition}

\begin{proof}
For every nonnegative integer \(u\),
\[
 (u)_j=(u)_2(u-2)_{j-2}
 \le (u)_2u^{j-2}
 \le (u)_2M_X^{j-2}
\]
whenever \(u=K_X(m)\).  Sum over \(m<X^2\), and apply
Proposition~\ref{prop:hm2}.
\end{proof}

\begin{hypothesis}[Strong anti-tail bound \((\mathrm{AT}'')\)]
\label{hyp:anti-tail-strong}
Uniformly over dyadic \(X\),
\[
 M_X\ll \lambda_X X^{o(1)}.
\]
Equivalently, for every \(\varepsilon>0\) there are
\(C_\varepsilon,X_\varepsilon>0\) such that
\[
 M_X\le C_\varepsilon X^\varepsilon\lambda_X
\]
for every dyadic \(X\ge X_\varepsilon\).  If \(\lambda_X=0\),
then \(M_X=0\), so this convention causes no ambiguity.
\end{hypothesis}

\begin{corollary}[\((\mathrm{AT}'')\) collapses the moment hierarchy]
\label{cor:anti-tail-full}
Assume Hypothesis~\ref{hyp:anti-tail-strong}.  Then
\((\mathrm{HM})_j\) holds for every fixed \(j\ge2\).  In particular,
\[
 G_n=e^{o(n)}
\qquad(n\to\infty)
\]
for every positive integer \(n\).
\end{corollary}

\begin{proof}
For fixed \(j\), Proposition~\ref{prop:moment-collapse} gives
\[
 \sum_{m<X^2}(K_X(m))_j
 \ll
 X^2\lambda_X^2
 \bigl(\lambda_XX^{o(1)}\bigr)^{j-2}
 =X^{2+o(1)}\lambda_X^j.
\]
This is precisely \((\mathrm{HM})_j\).  Taking any fixed
\(j>6\), for example \(j=8\), and applying
Theorem~\ref{thm:hm-pointwise} gives
\[
 \log G_n
 \ll n^{2/3+2/j+o(1)}
 =o(n).
\]
\end{proof}

\begin{remark}
Hypothesis~\((\mathrm{AT}'')\) is substantially stronger than any
fixed high-moment estimate, but it has the advantage that the complete
factorial-moment hierarchy follows from the already proved case
\(j=2\) in one line.  No independence assumption is used in
Proposition~\ref{prop:moment-collapse}.
\end{remark}


\section{An honest conditional reduction for the third moment}
\label{sec:honest-hm3-reduction}

We next isolate the two distinct inputs required by the complete
\(E_2\)-dispersion strategy.  Set
\[
 \Sigma_j(X)=\sum_{0\le m<X^2}(K_X(m))_j.
\]
For \(k,k'\in\mathbb F_p^\times\), define the multiplicative
correlation
\[
 M_p(k,k')
 =\#\{(r,r')\in\mathcal Z_p^2:kr=k'r'\}.
\]
It is the line-section statistic of the incidence cloud
\[
 \mathcal C_p
 =\{(r,h)\in\mathbb F_p^2:
       r\in\mathcal Z_p,\ r+h\in\mathcal Z_p\}.
\]
Indeed, writing \(r'=r+h\) gives the exact dictionary
\begin{equation}
 M_p(k,k')
 =\#\{(r,h)\in\mathcal C_p:
        k'h=(k-k')r\}.
 \label{eq:new-line-section-dictionary}
\end{equation}
The same kernel appears when a complete multiplicative square is
opened:
\begin{equation}
 \sum_{u\in\mathbb F_p^\times}
 F_p(ku)\overline{F_p(k'u)}
 =pM_p(k,k')-Z(p)^2.
 \label{eq:new-complete-multiplicative}
\end{equation}

For orientation, the scalar random-scale estimate suggested by
\eqref{eq:new-line-section-dictionary} is
\begin{equation}
 \sum_{\substack{k,k'\in\mathscr K\\k\ne k'}}
 M_p(k,k')
 \ll p^{o(1)}
 \left(
   |\mathscr K|Z(p)
   +\frac{|\mathscr K|^2Z(p)^2}{p}
 \right)
 \label{eq:new-mc-scalar}
\end{equation}
for short signed intervals \(\mathscr K\subset\mathbb F_p^\times\).
The complete \(E_2\)-dispersion argument requires a weighted,
operator-level version of \eqref{eq:new-mc-scalar}, because its
coefficients depend simultaneously on the Fourier frequency and on
the two auxiliary primes.  We record precisely the output required
from that weighted estimate.

\begin{hypothesis}[Weighted multiplicative CED estimate \((\mathrm{MC})\)]
\label{hyp:mc-ced}
Uniformly over dyadic \(X\),
\begin{equation}
 \Sigma_3(X)
 \ll
 X^{2+o(1)}\lambda_X^3
 +X^{-2/3+o(1)}M_X\Sigma_2(X).
 \label{eq:new-mc-ced}
\end{equation}
Here \((\mathrm{MC})\) means that the off-diagonal weighted form
arising from complete \(E_2\)-dispersion, whose complete local kernel
is \(pM_p(k,k')-Z(p)^2\) by
\eqref{eq:new-complete-multiplicative}, satisfies
\eqref{eq:new-mc-ced}.  Thus this is an operator-level
multiplicative-correlation hypothesis on the actual Ap\'ery zero sets,
not merely the unweighted scalar estimate
\eqref{eq:new-mc-scalar}.
\end{hypothesis}

\begin{hypothesis}[Anti-tail bound \((\mathrm{AT})\)]
\label{hyp:anti-tail-two-thirds}
Uniformly over dyadic \(X\),
\[
 M_X\ll X^{2/3+o(1)}\lambda_X.
\]
\end{hypothesis}

\begin{theorem}[\((\mathrm{MC})+(\mathrm{AT})\Rightarrow
 \mathrm{(HM)}_3\)]
\label{thm:mc-at-implies-hm3}
Hypotheses~\ref{hyp:mc-ced} and
\ref{hyp:anti-tail-two-thirds} imply
\[
 \Sigma_3(X)\ll X^{2+o(1)}\lambda_X^3.
\]
In other words, they imply \((\mathrm{HM})_3\).
\end{theorem}

\begin{proof}
By Proposition~\ref{prop:hm2},
\(\Sigma_2(X)\le5X^2\lambda_X^2\).  Substitution into
\eqref{eq:new-mc-ced}, followed by
Hypothesis~\ref{hyp:anti-tail-two-thirds}, gives
\begin{align*}
 \Sigma_3(X)
 &\ll X^{2+o(1)}\lambda_X^3
   +X^{-2/3+o(1)}
      \bigl(X^{2/3+o(1)}\lambda_X\bigr)
      \bigl(5X^2\lambda_X^2\bigr)\\
 &\ll X^{2+o(1)}\lambda_X^3.
\end{align*}
\end{proof}

\subsection{The complete \(E_2\)-dispersion architecture}

We explain why the multiplicative kernel
\eqref{eq:new-complete-multiplicative} is the natural one.  Let
\(L=X^2\), let \(p,q,\ell\in\mathcal P_X\) be distinct, and set
\(n=pq\ell\).  With
\[
 D_L(\alpha)=\sum_{0\le m<L}e(m\alpha),
\]
finite Fourier inversion gives
\begin{align}
 &\#\{0\le m<L:
       m\bmod p\in\mathcal Z_p,\
       m\bmod q\in\mathcal Z_q,\
       m\bmod\ell\in\mathcal Z_\ell\}\notag\\
 &\quad=
 \frac1{pq\ell}
 \sum_{a\bmod p}\sum_{b\bmod q}\sum_{c\bmod\ell}
 F_p(a)F_q(b)F_\ell(c)
 D_L\!\left(-\frac ap-\frac bq-\frac c\ell\right).
 \label{eq:new-trilinear-fourier}
\end{align}
The zero-frequency contribution, summed over ordered distinct prime
triples, is at most \(X^2\lambda_X^3\).

If \(m=m(r,s,t)\) is the CRT lift of
\((r\bmod p,s\bmod q,t\bmod\ell)\), then
\begin{equation}
 e(km/n)
 =
 e_p\!\left(k(r-s)(q\ell)^{-1}\right)
 e_\ell\!\left(k(t-s)(pq)^{-1}\right)
 e(ks/n).
 \label{eq:new-two-flip}
\end{equation}
Indeed, the integer
\[
 s+(r-s)q\ell(q\ell)^{-1}_{p}
   +(t-s)pq(pq)^{-1}_{\ell}
\]
has reductions \(r,s,t\) modulo \(p,q,\ell\), respectively.
Formula~\eqref{eq:new-two-flip} is the two-flip reciprocity step.

One now fixes the output prime \(p\), groups the other two primes
through the \(E_2\)-variable \(v=q\ell\), and applies Linnik's
dispersion method.  Opening the resulting square produces the
complete local correlation
\eqref{eq:new-complete-multiplicative}.  The diagonal frequencies are
controlled by Lemma~\ref{lem:fourier-nonconcentration}; the
off-diagonal frequencies are line sections
\eqref{eq:new-line-section-dictionary}.  Hypothesis~\((\mathrm{MC})\)
is exactly the weighted estimate needed to save the factor
\(X^{2/3}\) in this off-diagonal form, while
Hypothesis~\((\mathrm{AT})\) prevents the \(E_2\)-weights from
concentrating on a small exceptional set.  This is the same
diagonal/off-diagonal architecture as in Linnik dispersion and its
bilinear trace-function refinements; see
\cite{FouvryKM2014,KMS2017,Drappeau2017}.

\begin{remark}[Why the operator-level formulation is necessary]
\label{rem:new-ced-honesty}
The variable \(v=q\ell\) ranges over a sparse set of semiprimes, not
over a complete residue system, and the coefficient multiplying
\(F_p(k\bar v)\) depends on both \(v\) and \(k\).  Thus interval length
alone does not justify replacing the \(v\)-sum by
\eqref{eq:new-complete-multiplicative}, and a scalar Cauchy--Schwarz
argument cannot pull the remaining \(q,\ell\)-factor outside the
\(k\)-sum.  These are precisely the two losses absorbed into the
weighted hypothesis \((\mathrm{MC})\).  In particular,
\eqref{eq:new-mc-scalar} by itself is not claimed here to imply
\((\mathrm{HM})_3\).
\end{remark}

\subsection{The anchored-star obstruction}

\begin{proposition}[Anchored reflected star]
\label{prop:new-anchored-star}
For every sufficiently large integral \(X\), there is a family of
reflected, nonadjacent doublets
\(\mathcal Z_p^*\subset\mathbb F_p\), \(p\in\mathcal P_X\), such that
\[
 \lambda_X^*\asymp\frac1{\log X},
 \qquad
 M_X^*\asymp\frac{X}{\log X},
\]
and
\[
 \sum_{m<X^2}(K_X^*(m))_3
 \gg\frac{X^3}{(\log X)^3}.
\]
Consequently this family violates both \((\mathrm{AT})\) and
\((\mathrm{HM})_3\), and hence also violates the weighted conclusion
\((\mathrm{MC})\).
\end{proposition}

\begin{proof}
Choose \(m_0=\lfloor3X^2/4\rfloor\).  Apart from \(O(1)\) primes
dividing
\[
 (m_0-1)m_0(m_0+1)(m_0+2)(2m_0+1),
\]
put
\[
 \mathcal Z_p^*
 =\{m_0\bmod p,\ p-1-(m_0\bmod p)\},
\]
and use the empty set at the excluded primes.  The exclusions make
each nonempty set a reflected, nonadjacent doublet avoiding the
endpoint residues.  Every active prime marks \(m_0\), and the prime
number theorem and Mertens' theorem give
\[
 K_X^*(m_0)\asymp\frac X{\log X},
 \qquad
 \lambda_X^*=2\sum_{X<p\le2X}\frac1p+O(X^{-1})
 \asymp\frac1{\log X}.
\]
Therefore
\[
 \sum_{m<X^2}(K_X^*(m))_3
 \ge(K_X^*(m_0))_3
 \gg\frac{X^3}{(\log X)^3}.
\]
This exceeds \(X^{2+o(1)}(\lambda_X^*)^3\) by
\(X^{1-o(1)}\), while
\[
 \frac{M_X^*}
 {X^{2/3+o(1)}\lambda_X^*}
 =X^{1/3-o(1)}.
\]
Moreover, Proposition~\ref{prop:hm2} gives
\[
 \Sigma_2^*(X)\ll\frac{X^2}{(\log X)^2}.
\]
Hence the right-hand side of \eqref{eq:new-mc-ced}, evaluated on this
family, is at most
\[
 \frac{X^{2+o(1)}}{(\log X)^3}
 +X^{-2/3+o(1)}
   \frac{X}{\log X}\frac{X^2}{(\log X)^2}
 \ll\frac{X^{7/3+o(1)}}{(\log X)^3},
\]
which is \(o(X^3/(\log X)^3)\).  Thus the family also violates
\((\mathrm{MC})\).
\end{proof}

\begin{remark}
The anchored star satisfies reflection symmetry, has uniformly bounded
column size, and obeys the universal second-factorial-moment estimate.
Thus neither vertical zero-count bounds nor local palindromic
line-section information can replace the horizontal weighted content
of \((\mathrm{MC})\) and \((\mathrm{AT})\).
\end{remark}


\section{Computational evidence for the anti-tail hypotheses}
\label{sec:anti-tail-computation}

For each \(X=2^j\), \(4\le j\le12\), the following computation
enumerated every integer \(0\le m<X^2\) and every prime
\(X<p\le2X\).  The zero sets were generated by the division-free
recurrence for \((n!)^3b_n\).  The column \(m_X^{\mathrm{Pois}}\)
is the median of the maximum of \(X^2\) independent
\(\operatorname{Poisson}(\lambda_X)\) variables; explicitly, it is
the least integer \(t\) such that
\[
 \Pr\{\operatorname{Poisson}(\lambda_X)\le t\}^{X^2}\ge\frac12.
\]

\begin{center}
\begin{tabular}{rrrrrrrr}
\toprule
\(X\) & \(\lambda_X\) & \(M_X\) &
\(X^{2/3}\lambda_X\) &
\(\dfrac{M_X}{X^{2/3}\lambda_X}\) &
\(X\lambda_X\) &
\(\dfrac{M_X}{X\lambda_X}\) &
\(m_X^{\mathrm{Pois}}\)\\
\midrule
16   & 0.287426 & 2 & 1.825  & 1.096 & 4.599   & 0.435  & 3\\
32   & 0.169520 & 2 & 1.709  & 1.171 & 5.425   & 0.369  & 3\\
64   & 0.048016 & 2 & 0.768  & 2.603 & 3.073   & 0.651  & 2\\
128  & 0.157665 & 3 & 4.004  & 0.749 & 20.181  & 0.149  & 3\\
256  & 0.126286 & 3 & 5.092  & 0.589 & 32.329  & 0.0928 & 3\\
512  & 0.092772 & 3 & 5.937  & 0.505 & 47.499  & 0.0632 & 4\\
1024 & 0.082244 & 4 & 8.355  & 0.479 & 84.218  & 0.0475 & 4\\
2048 & 0.077521 & 4 & 12.502 & 0.320 & 158.763 & 0.0252 & 4\\
4096 & 0.076828 & 4 & 19.668 & 0.203 & 314.689 & 0.0127 & 4\\
\bottomrule
\end{tabular}
\end{center}

For the Poisson comparison, if \(Y_1,\ldots,Y_{X^2}\) are independent
\(\operatorname{Poisson}(\lambda_X)\) variables, the extreme-value
threshold \(t_X\) is determined by
\[
 X^2\Pr(Y_1\ge t_X)\asymp1.
\]
When \(\lambda_X\) stays bounded away from zero this gives
\[
 t_X\sim\frac{2\log X}{\log\log X};
\]
when \(\lambda_X\asymp1/\log X\), as suggested by the present data,
the same calculation gives
\[
 t_X\asymp\frac{\log X}{\log\log X}.
\]
Both scales are \(X^{o(1)}\), and therefore are compatible with
\((\mathrm{AT}'')\).  At the largest tested block the observed maximum
is \(M_X=4\), compared with the Poisson median \(4\);
it is \(4.92\) times smaller than \(X^{2/3}\lambda_X\) and
\(78.7\) times smaller than \(X\lambda_X\).

\begin{remark}
These finite computations support both \((\mathrm{AT})\) and the
stronger \((\mathrm{AT}'')\), but they do not prove either asymptotic
hypothesis.  Their main structural message is that the observed maximum
behaves like a sparse-occupancy extreme value rather than a power of
\(X\).
\end{remark}

% Required bibliography entry for proof.tex:
%
% \bibitem{Straub2014}
% A.~Straub,
% \emph{Multivariate Ap\'ery numbers and supercongruences of rational
% functions},
% Algebra Number Theory \textbf{8} (2014), no.~8, 1985--2008.
